\subsubsection{通信交接约束的恢复解释}
对关键运输任务$i$，若前段提供方可持续至相对时刻$\tau_i^{sw}$，后继中继$q$在绝对时刻$a_q$可用，含保护间隔$g$的交接要求可写为
\begin{equation}
 s_i+\tau_i^{sw}\ge a_q+g
 \quad\Longleftrightarrow\quad
 s_i\ge a_q-\tau_i^{sw}+g.
 \label{eq:q3_handoff_start}
\end{equation}
后继中继提前和前段保障延长，均可降低通信给任务开始施加的下界。当前同模型改进把这两方面结合，解释了约87.998 s工期下降。改进后两项运输末任务同刻返回，也说明下一轮位置调整可能遭遇新的资源瓶颈，收益不会无限线性叠加。

名义工期优先、增强通信保障和扰动后修复构成三类执行选择。附加0.25 dB保障备选比主方案增加约0.50\%工期、0.13\%能源；单项上线延迟则优先保持位置与路线、重新安排服务和运输时刻。12份修复给出了所测5、30、60 s上线延迟下的实际恢复方案，各自检查物理、时限、资源和连续通信。

\begin{table}[htbp]\centering\small
\caption{不同通信变化对应的计算与调整层次}\label{tab:q3_recovery_policy}
\begin{tabularx}{\linewidth}{p{2.3cm}LLL}\toprule
情景 & 处理方法 & 主要固定对象 & 交付结果\\\midrule
名义执行 & 使用工期优先方案 & 已验证任务、位置与资源结构 & 基准完整日程\\
持续附加损耗 & 对应保障备选与门限核验 & 所选备选的位置和任务 & 指定损耗条件下的日程\\
中继上线延后 & 连续时间修复 & 路线、位置、提供方与资源顺序 & 修复日程和额外代价\\
原覆盖关系失效 & 重新选择候选与任务 & 统一物理规则及实际需求 & 新结构经完整核验后使用\\\bottomrule
\end{tabularx}\end{table}

前三类已有保存方案或压力试验支持；覆盖结构重选是进一步扩展的求解入口。将扰动识别与调整权限明确对应，可以优先采用改动范围最小的恢复方式，同时在需要时返回上一决策层。通信保障由此体现为可选择的执行方案，而不止一个名义验证通过的标签。
